% Pursuing Stacks (cotes 134-2 to 134-8): archive cote 134-3, site page 34 (PDF page 35), typescript page 34.
% typescript. Draft compared: ps2.tex:1621-1693.
% AI-transcribed by Claude Opus 5.5 (claude-opus-5-5, adaptive thinking, effort high), Message Batches API, batch ps1 (msgbatch_01Rv5hCaNTtUt4JHx1yqoHPL), 2026-11-12; image review recorded in records/CHECK_DECISIONS_134_3_4.tsv.
\scanpage{134-3/34}
\mspage{\hand{34}}
\[ f\colon A \hand{\longrightarrow} A' \]
between small categories (or categories which are essentially small, name-%
ly equivalent to small categories) is called a weak equivalence, iff the
corresponding morphism of topoi
\[ f\hand{\hat{}}\colon \deleted{\illegible{1}}\ A\hand{\hat{}} \hand{\to} A'\hand{\hat{}} \]
is a weak equivalence.

Coming back to test categories A, which allow to construct the cor-%
responding modelizers $A\hand{\hat{}}$, our point of view here is rather that \uline{the test
categories are each just as good \inserted{one} as the other\deleted{s}}, and \hand{$\Delta$} just as good as
\hand{$\Box$} or any of the $T_{p,q}$ and not any better ! Maybe it's the one though which
turns out the most economical for computational use, the nerve functor
(Cat)$\hand{\to}$ \hand{$\Delta$}$\hand{\hat{}}$ being still the neatest known of all \inserted[typed]{model preserving} embeddings of (Cat) into
categories $A\hand{\hat{}}$ defined by modelizers.\hand{$^{(31)}$.} Another \inserted[typed]{point}, still more important \deleted{\xout{point}}
it seems to me, is that the natural functor
\[ \text{(Topoi)} \hand{\longrightarrow} \text{Pro(Hot)} \]
defined by the \hand{\v{}}Cech-Cartier-Verdier process, and which allows for another
description of weak equivalences of topoi\hand{,} \inserted{namely} as those made invertible by this
functor, are directly defined via semi-simplicial structures or simplicial
structures (of the type ``nerve of a covering'')\hand{$^{(32)}$.}

Modelizers of the type $A\hand{\hat{}}$, with A a test category, surely deserve
a name - lets call them \uline{elementary} \uline{modelizers}, as they correspond to \xout{xx}
the case of an ``elementary'' or ``trivial'' algebraic structure species,
whose set-theoretical realizations can be expressed as just \uline{any} functors
\[ A^{o} \hand{\longrightarrow} \text{(Sets)} , \]
without any exactness condition of any kind\over{,}{;} in other words they can
be viewed as just diagrams of sets of a specified type, with specified
commutativity relations. A somewhat ambitious question maybe is whether
on such a category \inserted[typed]{$M=$} $A\hand{\hat{}}$, namely an elementary modelizer, there \inserted[typed]{cannot} \xout{xxx} exist
\inserted[typed]{any} \xout{xx}other modelizing structure. In any case, the one we got is intrinsicly
determined in terms of M, which is a topos, by the prescription that
an arrow $f\colon a \hand{\to} b$ within M is a weak equivalence if and only if the
corresponding morphism \inserted[typed]{for} \xout{xx} the induced topoi $M_{/a}$ and $M/_{b}$ is a weak equiva-%
lence (in terms of the Artin-Mazur criterion \deleted{\unsure{above section}} \inserted{of section 17}\hand{)}.

\xout{Let} A more crucial question I feel is whether the \inserted[typed]{right} adjoint functor
$j_{A}$ to $i_{A}$ in (**) (cf. \deleted{\illegible{1}} \inserted{section 28)} is equally model preserving, whenever A is a
test category. This, as we have seen, is \uline{not} automatic\unsure{,} whenever we
have a model preserving functor between modelizers, \inserted{whenever} this functor admit\over{ting}{s}
an adjoint functor. In \inserted[typed]{more} \inserted{general} \xout{xxxxxxx} terms \xout{and more} still, let
% Hand-drawn arrows are marked blue; typed objects remain plain.
\diagram{\begin{tikzcd} M \ar[color=blue, r, shift left, "\hand{i}"] \& M' \ar[color=blue, l, shift left, "\hand{j}"] \end{tikzcd}}
\margin{\diagram{\begin{tikzpicture}[baseline=(n.base)]\node[draw,ellipse,inner sep=3pt] (n) {with a};\end{tikzpicture}}}
be a pair of adjoint functors, with M,M' endowed \deleted{with two} \inserted{with a} \uline{``saturated''} \deleted{sets}
